\documentclass[10pt]{article}
\usepackage[top=.25cm, bottom=1.5cm, left=1.6cm, right=1.6cm]{geometry}
\usepackage{amssymb,amsmath,setspace}
\title{TP Algorithmique et optimisation discrète}
\author{Ilès Ouyahia, Nicolas Montauriol}
\date{}
\renewcommand{\thesubsection}{Question~\arabic{subsection}}
\makeatletter
\@addtoreset{subsection}{section}
\makeatother
\begin{document}
\maketitle
\vspace{-.75cm}
\section{Équation de Bellman pour un paragraphe}

\subsection{}
Soit $D=[L_1,\ldots,L_\ell]$ un découpage optimal du paragraphe $P$. On note $P'$ le sous-paragraphe qui commence au premier mot de $L_2$, et $D'=[L_2,\ldots,L_\ell]$ le découpage correspondant. Par définition, le coût de ce découpage est de
\[
C(D)=N(e_1)+C(D')
\tag{1}
\]
Supposons, par l'absurde, que $D'$ ne soit pas optimal pour $P'$. Il existe alors un découpage optimal $D''$ de $P'$ tel que
\[
C(D'')<C(D')
\tag{2}
\]
Considérons maintenant le découpage $D'''$ de $P$ obtenu par concaténation de $L_1$ et $D''$ tel que $D''' = L_1 \cup D''$ \\
Par définition toujours ($L_1$ étant une ligne interne), le coût de ce découpage est de
\[
C(D''') \;=\; N(e_1)+C(D'')
\tag{3}
\]
\begin{align*}
&\text{d'après (2)} && N(e_1) + C(D'') \;<\; N(e_1) + C(D') \\
&\text{on remplace dans (3)} && C(D''') \;<\; N(e_1) + C(D') \\
&\text{on remplace encore, d'après (1)} && C(D''') \;<\; C(D)
\end{align*}
Ce résultat est en contradiction avec l'hypothèse d'optimalité de $D$. $D'$ est donc bien un découpage optimal de $P'$.

\subsection{}
Pour un paragraphe de \(n\) mots, avec \(r\) mots restants, et \(j\) mots par ligne, une équation de Bellman adaptée peut s'écrire
\[
\begin{aligned}
J(r) &=
\begin{cases}
0 & \text{si } r=0,\\[4pt]
\displaystyle \min_{1 \le j \le r}\Big( \mathrm{cost}(n-r+1,\,n-r+j) + J(r-j) \Big) & \text{sinon.}
\end{cases}\\[8pt]
\mathrm{cost}(i,k) &=
\begin{cases}
\infty & \text{si } \Delta(i,k) > M,\\
0 & \text{si } k = n,\\
\big(M - \Delta(i,k)\big)^3 & \text{sinon.}
\end{cases}
\end{aligned}
\]

\subsection{}
\subsubsection*{Comparaisons}
Chaque fonction \(J(r)\) n’est évaluée qu’une seule fois grâce à la mémoïsation, et chaque appel teste toutes les combinaisons de lignes possibles. Si l’on note \(t\) le nombre maximal de mots contenus dans une ligne de taille \(M\) (\(\mu = 1\), mots d’un seul caractère), on peut borner \(t\) par \(\left\lfloor \tfrac{M + 1}{2} \right\rfloor\). Le nombre de comparaisons à chaque appel de \(J\) est alors noté \(C(r)\) tel que

\[
C(r) = \min\!\Big(r,\,\max\{\,k - i + 1 \;\mid\; \Delta(i,k) \le M\,\}\Big)
       = \min(r,\,t)
\]
On en déduit le nombre total de comparaisons \(T\) tel que
\[
T = \sum_{r=1}^{n} C(r)
   = \sum_{r=1}^{n} \min(r,\,t)
   = \sum_{r=1}^{t} r \;+\; \sum_{r=t+1}^{n} t
   = \frac{t(t+1)}{2} + t(n-t)
   = t\,n - \frac{t(t-1)}{2}
\]
Puisque \(t\) ne dépend que de \(M\), on a alors \(T = \Theta(nM)\). On peut supposer que la taille d'une ligne  (\(M\)) est constante, ainsi \(t\) l’est aussi. Tant que \(t = O(1)\) (et donc \(\ll n\)), la complexité temporelle est linéaire :
\(
\boxed{T = \Theta(n)}
\)

\subsubsection*{Place mémoire}
La mémoïsation nécessite de stocker tous les \(J(r)\) déjà calculés dans un tableau de taille \(n\). Chaque case mémorise le coût minimal (ainsi qu'éventuellement l’indice du choix optimal), 
ce qui induit une complexité linéaire :
\(
\boxed{Memoire = \Theta(n)}
\)
\end{document}
